\rulesection{20.0 Optional Rules}

\rulesubsection{20.1 Solitaire Play}

The game can be played by one player moving both sides. The main obstacle is choosing combat options.

To resolve this, first choose an option for the attacker. Then roll one die. If the senior defending leader has a combat bonus of 2 or 3, subtract 1 from the die roll. Choose the defender's option as follows:
\begin{itemize}
\item 1--2: \textit{Pitched Battle} (if the defender is unsupplied, roll again)
\item 3--4: \textit{Stand}
\item 5--6: \textit{Withdraw}
\end{itemize}

\rulesubsection{20.2 Surprise Attacks}

A player may attempt a surprise attack if he force marches units into an enemy-occupied hex and all enemy units in the hex either fail to react or do not attempt reaction. Roll one die. If the result is less than the initiative of the senior moving leader, the player may make a surprise attack.

\textbf{20.21} The surprise attack is resolved immediately during the Movement Phase. Players do not choose combat options. The combat lasts two rounds unless one side is demoralized in the first. The defender applies none of the usual commitment modifiers. The attacker shifts one column right on the Combat Results Table. The defender may not fire during the first round.

\textbf{20.22} Units that have made a surprise attack may not move for the remainder of the phase.

\textbf{20.23} If both sides have units in the hex at the start of the Combat Phase, combat may occur there as usual.
